\documentclass[11pt]{article}

\usepackage[margin=1in]{geometry}
\usepackage{amsmath,amssymb}
\usepackage{booktabs}
\usepackage{enumitem}
\usepackage{xcolor}
\usepackage{hyperref}

\hypersetup{
  colorlinks=true,
  linkcolor=blue!60!black,
  citecolor=blue!60!black,
  urlcolor=blue!60!black
}

\setlist[itemize]{leftmargin=1.4em}
\setlist[enumerate]{leftmargin=1.6em}

\title{Improving the GQE Compiler for Unitary Synthesis:\\
Findings from the April 30, 2026 Code and Literature Review}
\author{Prepared for the \texttt{gqe-torch} project}
\date{April 30, 2026}

\begin{document}
\maketitle

\begin{abstract}
The 4-qubit benchmark logs show that the current GQE compiler can sometimes
recover high-fidelity circuits only after expensive post-training angle
refinement, but it is not yet learning compact high-fidelity circuits during the
policy phase.  During training, the Pareto best fidelity often remains around
\(0.55\)--\(0.58\), \texttt{best\_cnot@0.99} remains \(-1\), and the mean sampled
length often saturates the rollout cap.  This is consistent with a reward and
credit-assignment failure: structural penalties are disabled until fidelity
crosses \(0.99\), while the policy almost never reaches that region before the
final 500-step refinement.  This report reviews the relevant implementation,
connects it to recent synthesis literature, and proposes a staged improvement
plan.  Three important implementation fixes are now in place: archive entries
are verified from their exact stored token/angle pair, \texttt{STOP} has been
separated from inactive padding through a dedicated \texttt{PAD} token, and the
policy now treats angle refinement as a decoder with supervised angle-head
training.  The remaining highest-priority work is to make compactness visible
well before \(F=0.99\), add generator-aware baselines for brickwork targets,
reduce redundant refinement work, and introduce multi-target pretraining so
target conditioning becomes meaningful.
\end{abstract}

\section{Scope and Relationship to the Earlier Review}

This document updates and extends
\texttt{docs/unitary\_synthesis\_literature\_code\_review\_20260430.tex}.  That
earlier review contains useful background, but several implementation concerns
listed there are now partly or fully addressed in the current code:
\begin{itemize}
  \item The rollout context length is treated as an engineering cap, and
  rollout exits once all batch rows emit \texttt{STOP}.
  \item The policy includes target conditioning, wrapped angle handling, and a
  proper mask for non-parametric angle supervision.
  \item The Pareto archive now includes total gate count in dominance checks.
  \item Archive refinement now re-evaluates the exact stored token/angle pair
  before insertion into the rebuilt archive.
  \item \texttt{STOP} is reserved for termination, with \texttt{PAD} used for
  inactive suffix positions.
  \item The circuit simplifier has been removed from the training and
  reporting pipeline.
  \item The generator is now trained as a skeleton policy: discrete PPO uses
  token likelihoods, while angle heads learn refined angles by supervised
  periodic regression.
  \item Benchmarking now records threshold-constrained metrics, verifies Qiskit
  fidelity, supports seed sweeps, and stores richer metadata.
\end{itemize}

The remaining problems are therefore more specific: the policy is being trained
on a weak compactness signal, the final-refinement budget is still spent on too
many redundant skeletons, target conditioning is weak in one-target runs, and
the benchmark baseline underuses known structure in the generated brickwork
targets.

\section{What the 4-Qubit Logs Say}

The logs from seeds 1020--1026 show a repeated pattern:
\begin{itemize}
  \item In-policy fidelity is low.  \texttt{F\_seen} and
  \texttt{F\_pareto\_best} are generally near \(0.55\)--\(0.58\), and
  \texttt{best\_cnot@0.99=-1} throughout training.
  \item Final fidelity is mostly produced by post-training refinement.  For
  example, completed runs report final refined GQE fidelities such as
  \(0.9922\), \(0.9945\), and \(0.9939\), but also failures or marginal cases
  such as \(0.9877\), \(0.9818\), \(0.9882\), and \(0.9502\).
  \item Sampled sequence length often saturates near the cap
  (\(\texttt{mean\_len}\approx 458\)).  This implies that the model has not
  learned termination as a useful decision, so both evaluation and refinement
  operate on bloated skeletons.
  \item The final circuits are not competitive on CNOT count.  Qiskit reports
  \(95\) CNOTs for these dense \texttt{UnitaryGate} baselines, while GQE often
  reports \(103\)--\(152\) CNOTs after refinement.  More importantly, for the
  current brickwork target generator, an oracle decomposition should be much
  smaller than either number, as explained in Section~\ref{sec:brickwork}.
\end{itemize}

This is not mainly an optimizer-step-count problem.  It is a search-space and
learning-signal problem: the policy is not being asked, soon enough or strongly
enough, to prefer short, low-CNOT skeletons that are still refinable.

\section{Implementation Findings}

\subsection{[DONE] High-priority correctness issue: archive refinement can report
inconsistent circuits}

The previous archive-refinement path could combine a fidelity from one circuit
with structure metrics from another.  That was the most urgent code-level risk
because it could directly inflate final reported GQE results.

The current implementation removes the circuit simplifier and rebuilds the
refined archive only from executable token/angle pairs.  For each archive entry,
\texttt{refine\_pareto\_archive} constructs the stored token row, initializes
from the stored angles, runs angle refinement, and then re-evaluates the exact
token/angle pair that will be inserted.  If refinement lowers fidelity, the
original stored angles are kept and the fidelity is recomputed from that
original pair.  Reporting therefore no longer combines old fidelity with new
structure.

Remaining guardrail: keep the regression test that forces a refiner to return a
worse angle vector and verifies that the archive stores the re-evaluated
original pair.  Any future canonicalization, pruning, or skeleton deduplication
must preserve the same invariant: \((F,\mathrm{CNOT},D,G)\) must correspond to
one stored token/angle pair.

\subsection{Structural pressure is mostly inactive during training}

The lexicographic reward in \texttt{trainer.py:571--585} is
\[
  C(F,S) =
  -\gamma^{L-1} R_F
  + \mathbf{1}\{F \ge F_{\mathrm{struct}}\}
  (w_C C_{\mathrm{CNOT}} + w_D D)
  + C_{\mathrm{nostop}},
  \qquad
  R_F = -w_F \log(1-F+\epsilon),
\]
with \(F_{\mathrm{struct}}=0.99\) in \texttt{config.yml}.  The 4-qubit
logs show that the policy phase almost never produces \(F \ge 0.99\).
Consequently the structure objective is inactive during the period when the
model learns its sequence distribution.  The model is rewarded primarily for
refinable fidelity, not for compactness.

This explains the observed saturation at long lengths and the weak CNOT results.
The final 500-step optimizer can tune angles, but it cannot turn a long, poorly
chosen skeleton into a minimal one.

Recommended replacement:
\begin{itemize}
  \item Use a fidelity curriculum: start structure pressure at a low threshold
  such as \(0.4\)--\(0.6\), then raise it as the archive improves.
  \item Or use smooth gated pressure,
  \[
    C = w_F \log(1-F+\epsilon)
      + \sigma(k(F-\tau)) (w_C C_{\mathrm{CNOT}} + w_D D + w_G G),
  \]
  where \(\tau\) is annealed upward.
  \item Maintain separate preference labels inside fidelity bands:
  among candidates with similar refined fidelity, prefer lower CNOT, depth, and
  total gates.
  \item Make the \texttt{STOP} decision part of the main objective, not only a
  terminal penalty for never stopping.
\end{itemize}

\subsection{Best raw rollout selection is misleading}

\texttt{trainer.py} still tracks the highest raw pre-refinement fidelity and
then spends 500 Adam steps refining that candidate after training.  The logs show
``best raw'' fidelities near \(0.04\)--\(0.06\), while the in-epoch refined
fidelity can be much higher.  This means the post-training ``best raw'' path is
not selecting the best skeleton; it is selecting the lucky raw-angle sample.

The final expensive refinement budget should instead be spent on:
\begin{itemize}
  \item archive entries chosen by refined fidelity and structural objectives,
  \item top-\(k\) skeletons under short inner refinement, with duplicate
  skeletons merged,
  \item diverse low-CNOT candidates, even if their raw angles are poor.
\end{itemize}

\subsection{[DONE] \texttt{STOP} is overloaded as both termination and no-op}

The old design used \texttt{STOP} both as sequence termination and as an
operational no-op.  That made length statistics ambiguous: components using
``first \texttt{STOP}'' saw a truncated circuit, while evaluators could still
interpret later \texttt{STOP} tokens as no-ops.

The current token layout reserves \texttt{STOP} for termination and adds a
dedicated \texttt{PAD} token for inactive suffix positions.  Rollout sampling
now emits one terminal \texttt{STOP} and then fills the rest of the fixed JAX
shape with \texttt{PAD}.  \texttt{PAD} is masked from active sampling, and both
\texttt{STOP} and \texttt{PAD} are treated as no-op tokens by structure and
unitary evaluation.  This makes sampled length, no-stop rate, and active-gate
statistics meaningful.

Remaining guardrail: future editing or canonicalization passes should preserve
the invariant ``active gates, one \texttt{STOP}, then \texttt{PAD} suffix'' and
must not introduce interior \texttt{STOP} tokens.

\subsection{Target conditioning is currently a constant bias for one-target
runs}

\texttt{trainer.py} flattens the target unitary into real and imaginary
features, and \texttt{policy.py} injects their projection as an additive
bias at every position.  For a single-target run, every sample in every epoch has
the same target features.  The target projection therefore behaves like a
learned constant offset rather than a meaningful conditional policy.

Target conditioning becomes useful only with multi-target training or
pretraining.  It should also be made invariant to global phase and better
structured than a flat matrix.  Candidate encoders include:
\begin{itemize}
  \item global-phase-canonicalized real/imaginary unitary features for small
  \(n\),
  \item residual-unitary features \(R_t = U_{\mathrm{target}} U_t^\dagger\),
  \item local two-qubit marginal or operator-Schmidt features for larger \(n\),
  \item problem-family metadata such as brickwork layer, coupling graph, and
  target fidelity/cost tradeoff token.
\end{itemize}

\subsection{A small runtime issue in rollout}

\texttt{policy.py} recomputes
\(\texttt{target\_features @ target\_proj + bias}\) inside every decode step.
For one target and fixed batch features, this projection can be computed once
outside the autoregressive loop and passed into \texttt{decode\_step}.  This will
not solve the 700-second per-seed runtime alone, but it is a clean reduction in
unnecessary matrix multiplication.

\section{Benchmark Diagnosis: Brickwork Targets Have a Hidden Oracle}
\label{sec:brickwork}

\texttt{target.py:82--105} generates \texttt{brickwork\_haar} targets by applying
alternating nearest-neighbor random two-qubit unitaries.  With four qubits and
the default depth \(2n=8\), the target contains twelve adjacent two-qubit Haar
blocks: two blocks on even layers and one block on odd layers.  Shende, Markov,
and Bullock, and independently Vidal and Dawson, show that an arbitrary
two-qubit unitary can be implemented with three CNOT gates plus single-qubit
rotations~\cite{shende2004minimal,vidal2004three}.  Therefore the generator
itself provides a constructive exact upper bound of roughly
\[
  12 \times 3 = 36
\]
CNOT gates before any cross-layer cancellation.

This makes the current comparison misleading.  Qiskit is given only a dense
\texttt{UnitaryGate} in \texttt{benchmark.py:127--147}, so it must rediscover
structure from a \(16 \times 16\) matrix.  Reporting Qiskit at \(95\) CNOTs is
useful as a dense-unitary baseline, but it is not a fair structured-target
baseline.  GQE at \(100\)--\(150\) CNOTs is also far from the generator-aware
oracle.

Recommended baselines:
\begin{itemize}
  \item \textbf{Generator oracle}: store the brickwork recipe and decompose each
  two-qubit block with KAK/three-CNOT synthesis.
  \item \textbf{Qiskit original-circuit baseline}: build the brickwork circuit
  explicitly, then transpile it, instead of first collapsing it to a dense
  \texttt{UnitaryGate}.
  \item \textbf{BQSKit/QSearch/LEAP/QFAST}: compare against synthesis tools
  designed for small-unitary and structured synthesis~\cite{bqskitdocs,qfast,leap}.
  \item \textbf{Qiskit AQC}: include approximate quantum compiling with fixed
  CNOT networks and L-BFGS-B optimization~\cite{madden2021aqc,qiskitaqc}.
\end{itemize}

The long-term GQE target should be not merely to beat the dense Qiskit baseline,
but to approach the generator-aware oracle when the target family has compact
latent structure.

\section{Literature Synthesis}

\subsection{Exact and numerical synthesis}

Classical unitary synthesis provides several lessons:
\begin{itemize}
  \item Exact decompositions give important scale checks.  Worst-case
  \(n\)-qubit dense-unitary lower bounds are exponentially large, while the
  current brickwork targets are not dense generic unitaries in that sense.
  Shende--Bullock--Markov top-down synthesis and two-qubit optimal synthesis
  should be treated as sanity baselines~\cite{shende2006synthesis,shende2004minimal}.
  \item BQSKit separates search, synthesis, and instantiation.  QSearch searches
  structures with numerical optimization; LEAP reduces candidate explosion by
  prefixing; QFAST uses hierarchical continuous blocks to reduce branching
  factor~\cite{bqskitdocs,qfast,leap}.  GQE is currently closer to brute-force
  token generation plus optimizer, and should borrow the hierarchical and
  prefixing ideas.
  \item Approximate quantum compiling formulates synthesis as fixed-template
  optimization with hardware constraints and can compress circuits with small
  loss in fidelity~\cite{madden2021aqc,robertson2023aqc}.  This is directly
  relevant to post-GQE refinement: choose a compact skeleton first, then perform
  strong continuous optimization.
  \item Recent work by Gomathi and Meiburg argues that simple gradient descent
  can reliably find optimal circuits for generic unitaries when the skeleton is
  adequately parameterized and avoids parameter-deficient patterns~\cite{gomathi2026gd}.
  This supports using carefully designed skeleton families instead of arbitrary
  random token streams.
\end{itemize}

\subsection{Machine-learning approaches}

The ML literature suggests that the current PPO-style online loop is not the
only useful training signal:
\begin{itemize}
  \item The original GQE paper uses transformer-style circuit generation and
  discusses pretraining for quantum simulation tasks~\cite{nakaji2024gqe}.
  This supports using offline corpora of good circuits rather than learning
  every target from scratch.
  \item Conditional-GQE shows the importance of context-aware generation and
  preference-based training over many related problem instances~\cite{minami2025conditional}.
  In this project, target conditioning will not pay off until training batches
  contain many targets.
  \item Persistent-DPO and hybrid offline/online training were proposed to
  improve GQE convergence and avoid limitations of ordinary DPO~\cite{nakamura2025pdpo}.
  The same idea maps naturally to circuit synthesis: keep a persistent archive
  of good skeletons and update the generator by preferences inside fidelity
  bands.
  \item Deep RL compilers demonstrate that architecture-aware rewards can reduce
  depth/gate count and can exploit hardware constraints~\cite{moro2021rl,foesel2021rl,wang2024rl}.
  GQE should include coupling graph and direction constraints early in the
  action space, not only after synthesis.
  \item Diffusion models for circuit synthesis show that conditional generative
  models can generate and edit circuits under constraints~\cite{furrutter2024diffusion}.
  This is a useful alternative for pretraining or repair, even if the main model
  remains autoregressive.
  \item GFlowNets are designed to sample diverse high-reward objects rather than
  collapse onto one mode~\cite{jain2022gflownet}.  A GFlowNet-style objective is
  attractive for Pareto synthesis because the archive needs diverse
  high-fidelity, low-cost circuits.
\end{itemize}

\section{Recommended Compiler Redesign}

\subsection{[DONE] Separate skeleton search from angle optimization}

The model still samples angle values during rollout, but the training objective
now treats the discrete token stream as the primary RL action.  The useful
score is the refined skeleton score produced by \texttt{AngleRefiner}, and the
angle head is trained separately toward refined rotation angles using a masked
periodic regression loss.

This implements the first stage of a more robust architecture:
\begin{enumerate}
  \item The generator is rewarded for choosing a refinable token skeleton:
  CNOT locations, fixed gates, and rotation axes.
  \item A continuous optimizer fits angles for the skeleton.
  \item The generator is updated from the refined skeleton score, while angle
  heads are trained by supervised regression to refined angles.
\end{enumerate}

This makes the credit assignment honest.  The discrete generator is rewarded for
choosing refinable low-cost skeletons; the continuous optimizer is responsible
for numerical fitting.

Remaining gap: the skeleton is still a raw gate-level sequence, not yet a
compact macro-block representation.  The next architectural step is to emit
larger skeleton actions such as local SU(4) blocks, hardware entangling layers,
or brickwork templates.

\subsection{Use macro actions}

For 4-qubit and 8-qubit synthesis, emitting raw one-gate tokens is too long a
horizon.  Add actions such as:
\begin{itemize}
  \item adjacent arbitrary two-qubit block, later instantiated by KAK,
  \item hardware-native entangling layer plus local Euler rotations,
  \item brickwork layer template,
  \item delete/replace/swap local window actions for circuit editing.
\end{itemize}

The raw gate-level circuit can still be produced at the end, but search should
operate over a much shorter and more expressive representation.

\subsection{Train with preference bands rather than one scalar reward}

Within each epoch, construct pairwise or listwise preferences:
\begin{itemize}
  \item If fidelities differ by more than a band width, prefer higher fidelity.
  \item If fidelities are in the same band, prefer lower CNOT count, then depth,
  then total gates, then shorter length.
  \item Give extra preference to candidates that pass moving thresholds such as
  \(0.6\), \(0.8\), \(0.9\), \(0.95\), and \(0.99\).
\end{itemize}

This mirrors DPO/CPO/Persistent-DPO ideas while directly encoding the
multi-objective goal.  It also avoids the present failure mode where structure
is invisible until \(F\ge 0.99\).

\subsection{Make termination learnable}

The model should see explicit positive examples of short circuits.  Recommended
changes:
\begin{itemize}
  \item Add a length prior or survival penalty at every non-\texttt{STOP} step.
  \item Use iterative deepening: train/search with small caps first, increase
  only when no candidate crosses the moving fidelity threshold.
  \item Merge duplicate skeletons before refinement so repeated long variants do
  not consume the budget.
  \item Report active gates up to the first \texttt{STOP}, not the fixed JAX
  context length.
\end{itemize}

\section{Runtime Improvement Plan}

The current logs are roughly 700 seconds per completed seed.  The fastest
improvements are not low-level micro-optimizations; they reduce the number and
length of expensive circuit evaluations.

\begin{enumerate}
  \item \textbf{Refine fewer candidates.}  Run the inner refiner only on
  top-\(k\), archive-eligible, or novelty-filtered skeletons.  Score the rest
  with raw or very short refinement.
  \item \textbf{Deduplicate skeletons.}  Hash the active token prefix up to
  \texttt{STOP} before refinement.  Keep the best angle initialization per
  skeleton.
  \item \textbf{Hoist target conditioning.}  Precompute the projected target
  bias outside the autoregressive decode loop.
  \item \textbf{Use active-length buckets.}  Evaluate/refine only active-token
  prefixes where possible rather than full context-length rows.  If JAX shape
  constraints make ragged batches awkward, bucket by length and pad within the
  bucket.
  \item \textbf{Use stronger optimizers on fewer skeletons.}  Multi-start
  L-BFGS-B or an AQC-style objective may outperform many cheap Adam refinements
  on poor skeletons.
  \item \textbf{Cache prefix/suffix products for sweeps.}  The evaluator scans
  all gates with \texttt{lax.scan}; local angle sweeps and window replacements
  can reuse prefix/suffix unitaries.
  \item \textbf{Prefer local/hardware skeletons.}  For nearest-neighbor
  brickwork targets, avoid sampling all-to-all directed CNOTs unless the target
  family or hardware requires them.
\end{enumerate}

\section{Pretraining Plan}

Pretraining is likely to help, but only if it is aligned with the synthesis
objective rather than just next-token imitation on arbitrary circuits.

\subsection{Data}

Build an offline corpus with multiple target families:
\begin{itemize}
  \item exact generator recipes for \texttt{random\_reachable} and
  \texttt{brickwork\_haar};
  \item KAK-decomposed two-qubit Haar blocks;
  \item Qiskit, BQSKit, QSearch, LEAP, QFAST, and AQC outputs for small
  unitaries;
  \item self-play archive entries from successful GQE runs, verified by
  re-evaluating the stored token/angle pair;
  \item perturbed and partially deleted circuits paired with their residual
  unitaries.
\end{itemize}

\subsection{Objectives}

Use several losses:
\begin{itemize}
  \item behavior cloning for compact skeletons and \texttt{STOP};
  \item denoising: recover a good next block from a residual unitary and partial
  circuit;
  \item preference training: choose lower-CNOT circuits within fidelity bands;
  \item angle warm-start regression toward refined angles, with periodic loss;
  \item multi-objective conditioning: target fidelity threshold, CNOT budget,
  coupling graph, and desired depth budget as input tokens/features.
\end{itemize}

\subsection{Batching}

Train on many targets per batch.  Otherwise target conditioning is just a
constant bias.  For small \(n\), flat unitary features are acceptable after
global-phase canonicalization.  For larger \(n\), switch to residual and local
features, or train a blockwise policy that never consumes a full \(2^n \times
2^n\) matrix.

\section{Concrete Roadmap}

\subsection{One-week fixes}

\begin{enumerate}
  \item Add a generator-aware brickwork baseline and report it next to dense
  Qiskit.
  \item Lower or smooth the structural reward threshold and log the fraction of
  samples where structural penalties are active.
  \item Deduplicate active skeleton prefixes before final archive refinement.
  \item Hoist target projection out of the rollout loop.
  \item Stop spending the full post-training budget on the highest raw-fidelity
  rollout when archive or short-refined skeleton candidates are available.
\end{enumerate}

\subsection{One-month changes}

\begin{enumerate}
  \item Replace scalar GRPO/PPO reward with fidelity-band preferences.
  \item Move from raw gate-level skeletons to macro actions such as local
  SU(4)/KAK blocks and hardware-native entangling layers.
  \item Add iterative-deepening caps and a length/termination curriculum.
  \item Run BQSKit/QSearch/LEAP/QFAST/AQC baselines for 3-qubit and 4-qubit
  targets.
  \item Train on many targets per batch so target conditioning becomes an
  actual conditional signal instead of a per-run constant bias.
\end{enumerate}

\subsection{Longer-term redesign}

\begin{enumerate}
  \item Move to macro-block generation: local SU(4)/KAK blocks and entangling
  layers first, raw gates second.
  \item Pretrain on multi-target corpora and fine-tune per target with
  preference optimization.
  \item Add circuit editing actions so the model can improve an existing
  skeleton instead of generating a full cap-length sequence from scratch.
  \item For 8-qubit targets, synthesize hierarchically by partitions and
  residual blocks rather than monolithic token streams.
\end{enumerate}

\section{Conclusion}

The logs are consistent with a compiler that has a useful angle-refinement
backend but a weak structure-learning frontend.  The archive consistency,
\texttt{STOP}/\texttt{PAD}, and skeleton-angle credit-assignment fixes remove
the most serious bookkeeping problems.  The most important remaining shift is
to make compactness visible before \(F=0.99\), compare against baselines that
respect the target generator's latent structure, and move toward a
macro-skeleton, preference-trained, target-conditioned compiler with strong
continuous optimization as a decoder.

\begin{thebibliography}{99}

\bibitem{shende2004minimal}
V. V. Shende, I. L. Markov, and S. S. Bullock.
\newblock Minimal universal two-qubit controlled-NOT-based circuits.
\newblock \emph{Physical Review A} 69, 062321, 2004.
\newblock \url{https://www.nist.gov/publications/minimal-universal-two-qubit-controlled-not-based-circuits}.

\bibitem{vidal2004three}
G. Vidal and C. M. Dawson.
\newblock Universal quantum circuit for two-qubit transformations with three
controlled-NOT gates.
\newblock \emph{Physical Review A} 69, 010301, 2004.
\newblock \url{https://www.osti.gov/biblio/20640567}.

\bibitem{shende2006synthesis}
V. V. Shende, S. S. Bullock, and I. L. Markov.
\newblock Synthesis of quantum logic circuits.
\newblock \emph{IEEE Transactions on Computer-Aided Design of Integrated
Circuits and Systems} 25(6), 2006.
\newblock \url{https://www.nist.gov/publications/synthesis-quantum-logic-circuits}.

\bibitem{qfast}
E. Younis, K. Sen, K. Yelick, and C. Iancu.
\newblock QFAST: Quantum Synthesis Using a Hierarchical Continuous Circuit
Space.
\newblock arXiv:2003.04462, 2020.
\newblock \url{https://arxiv.org/abs/2003.04462}.

\bibitem{bqskitdocs}
Berkeley Quantum Synthesis Toolkit documentation.
\newblock Quantum Synthesis Explained and BQSKit overview.
\newblock \url{https://bqskit.readthedocs.io/en/latest/intro/synthesis.html};
\url{https://bqskit.lbl.gov/}.

\bibitem{leap}
E. Smith, M. G. Davis, J. M. Larson, E. Younis, L. Bassman,
W. Lavrijsen, and C. Iancu.
\newblock LEAP: Scaling Numerical Optimization Based Synthesis Using an
Incremental Approach.
\newblock \emph{ACM Transactions on Quantum Computing}, 2022.
\newblock \url{https://doi.org/10.1145/3548693}.

\bibitem{madden2021aqc}
L. Madden and A. Simonetto.
\newblock Best Approximate Quantum Compiling Problems.
\newblock arXiv:2106.05649, 2021.
\newblock \url{https://arxiv.org/abs/2106.05649}.

\bibitem{qiskitaqc}
IBM Quantum documentation.
\newblock Qiskit Approximate Quantum Compiler.
\newblock \url{https://qiskit.qotlabs.org/docs/api/qiskit/qiskit.synthesis.unitary.aqc.AQC}.

\bibitem{robertson2023aqc}
N. F. Robertson, A. Akhriev, J. Vala, and S. Zhuk.
\newblock Approximate Quantum Compiling for Quantum Simulation: A Tensor Network
based approach.
\newblock arXiv:2301.08609; \emph{ACM Transactions on Quantum Computing}, 2025.
\newblock \url{https://arxiv.org/abs/2301.08609}.

\bibitem{qiskitunitary}
IBM Quantum documentation.
\newblock Qiskit \texttt{UnitarySynthesis} pass and synthesis plugins.
\newblock \url{https://docs.quantum.ibm.com/api/qiskit/qiskit.transpiler.passes.UnitarySynthesis}.

\bibitem{nakaji2024gqe}
K. Nakaji et al.
\newblock The generative quantum eigensolver (GQE) and its application for
ground state search.
\newblock arXiv:2401.09253, 2024; revised 2025.
\newblock \url{https://arxiv.org/abs/2401.09253}.

\bibitem{minami2025conditional}
S. Minami, K. Nakaji, Y. Suzuki, A. Aspuru-Guzik, and T. Kadowaki.
\newblock Generative quantum combinatorial optimization by means of a novel
conditional generative quantum eigensolver.
\newblock arXiv:2501.16986, 2025.
\newblock \url{https://arxiv.org/abs/2501.16986}.

\bibitem{nakamura2025pdpo}
J. Nakamura and S. Sanji.
\newblock Persistent-DPO: A novel loss function and hybrid learning for
generative quantum eigensolver.
\newblock arXiv:2509.08351, 2025.
\newblock \url{https://arxiv.org/abs/2509.08351}.

\bibitem{moro2021rl}
L. Moro, M. G. A. Paris, M. Restelli, and E. Prati.
\newblock Quantum compiling by deep reinforcement learning.
\newblock \emph{Communications Physics} 4, 178, 2021.
\newblock \url{https://www.nature.com/articles/s42005-021-00684-3}.

\bibitem{foesel2021rl}
T. Foesel, M. Y. Niu, F. Marquardt, and L. Li.
\newblock Quantum circuit optimization with deep reinforcement learning.
\newblock arXiv:2103.07585, 2021.
\newblock \url{https://arxiv.org/abs/2103.07585}.

\bibitem{wang2024rl}
Z. T. Wang et al.
\newblock Quantum Compiling with Reinforcement Learning on a Superconducting
Processor.
\newblock arXiv:2406.12195, 2024.
\newblock \url{https://arxiv.org/abs/2406.12195}.

\bibitem{rotosolve}
M. Ostaszewski, E. Grant, and M. Benedetti.
\newblock Structure optimization for parameterized quantum circuits.
\newblock \emph{Quantum} 5, 391, 2021.
\newblock \url{https://quantum-journal.org/papers/q-2021-01-28-391/}.

\bibitem{furrutter2024diffusion}
F. Fuerrutter, G. Munoz-Gil, and H. J. Briegel.
\newblock Quantum circuit synthesis with diffusion models.
\newblock \emph{Nature Machine Intelligence} 6, 515--524, 2024.
\newblock \url{https://www.nature.com/articles/s42256-024-00831-9}.

\bibitem{jain2022gflownet}
M. Jain et al.
\newblock Biological Sequence Design with GFlowNets.
\newblock \emph{Proceedings of the 39th International Conference on Machine
Learning}, PMLR 162:9786--9801, 2022.
\newblock \url{https://proceedings.mlr.press/v162/jain22a.html}.

\bibitem{gomathi2026gd}
J. Gomathi and A. Meiburg.
\newblock Gradient descent reliably finds depth- and gate-optimal circuits for
generic unitaries.
\newblock arXiv:2601.03123, 2026.
\newblock \url{https://arxiv.org/abs/2601.03123}.

\end{thebibliography}

\end{document}
